\documentclass[11pt]{article}

\usepackage[margin=1in]{geometry}
\usepackage{amsmath,amssymb,amsfonts,mathtools,bm}
\usepackage{booktabs,longtable,array}
\usepackage{hyperref}
\usepackage{enumitem}
\usepackage{setspace}
\usepackage{xcolor}
\usepackage{listings}
\usepackage{graphicx}
\usepackage{float}

\setstretch{1.12}
\hypersetup{colorlinks=true,linkcolor=blue,urlcolor=blue,citecolor=blue}

\lstset{
    basicstyle=\ttfamily\small,
    breaklines=true,
    frame=single,
    columns=fullflexible
}

\title{Detailed Description of the HOOMD-blue Patchy Rubisco--EPYC1 Brownian Dynamics Model}
\author{}
\date{}

\begin{document}

\maketitle
\tableofcontents
\newpage

\section{Purpose and scope}

This document describes the latest HOOMD-blue version of the Rubisco--EPYC1 patchy-particle Brownian dynamics model. It is intended to provide enough information to write the model from scratch, explain the model in a methods section, and understand how the current simulation outputs are generated.

The model is designed to capture the following biological and physical features:
\begin{enumerate}[label=\arabic*.]
    \item Rubisco acts as a multivalent scaffold.
    \item EPYC1 acts as a flexible multivalent linker.
    \item Rubisco has explicit patch sites that can be contacted by EPYC1 sticker beads.
    \item EPYC1 has multiple sticker beads distributed along a polymer chain.
    \item Assembly is driven by continuous force-based interactions, not irreversible bookkeeping.
    \item Dynamics are Brownian and run on GPU through HOOMD-blue.
    \item The workflow preserves all requested outputs: GSD trajectory, XYZ trajectory for OVITO, CSV time series, and JPEG plots.
\end{enumerate}

\section{Model overview}

The system contains two biological components:
\begin{itemize}
    \item Rubisco, represented as one central scaffold particle plus eight patch particles.
    \item EPYC1, represented as a bead-spring polymer with spacer beads and sticker beads.
\end{itemize}

The central interaction is between EPYC1 sticker beads, type \(T\), and Rubisco patch particles, type \(P\). This interaction is a continuous attractive potential. There are no discrete binding flags in the HOOMD dynamics layer. Instead, contacts are detected during analysis by measuring sticker--patch distances.

\section{Particle types}

The model uses four particle types:
\begin{longtable}{>{\ttfamily}p{0.15\textwidth} p{0.25\textwidth} p{0.50\textwidth}}
\toprule
Type & Meaning & Role \\
\midrule
R & Rubisco center & Central scaffold particle for one Rubisco complex. \\
P & Rubisco patch & Patch site tethered to a Rubisco center; attracts EPYC1 stickers. \\
S & EPYC1 spacer & Non-sticky bead in the EPYC1 polymer chain. \\
T & EPYC1 sticker & Sticky EPYC1 bead that attracts Rubisco patch particles. \\
\bottomrule
\end{longtable}

\section{Particle numbers and ratios}

The default model uses:
\[
N_R = 24
\]
Rubisco scaffolds and:
\[
N_C = 48
\]
EPYC1 chains.

Thus, the EPYC1:Rubisco ratio is:
\[
\frac{N_C}{N_R} = \frac{48}{24} = 2.
\]

Each Rubisco has eight patch particles:
\[
N_{\mathrm{patches/Rubisco}} = 8.
\]

The total number of patch particles is:
\[
N_P = 24 \times 8 = 192.
\]

Each EPYC1 chain has 12 beads:
\[
N_{\mathrm{beads/EPYC1}} = 12.
\]

The EPYC1 bead sequence is:
\[
[S,T,S,T,S,T,S,T,S,T,S,T].
\]

Each EPYC1 has six sticker beads:
\[
N_{\mathrm{stickers/EPYC1}} = 6.
\]

The total number of EPYC1 beads is:
\[
N_{\mathrm{EPYC1\ beads}} = 48 \times 12 = 576.
\]

The total number of EPYC1 sticker beads is:
\[
N_T = 48 \times 6 = 288.
\]

The total number of simulated particles is:
\[
N_{\mathrm{total}} = 24 + 192 + 576 = 792.
\]

The sticker-to-patch ratio is:
\[
\frac{N_T}{N_P} = \frac{288}{192} = 1.5.
\]

This ratio is important because it means there are more EPYC1 sticker sites than Rubisco patch sites. This promotes bridging and network formation, but also requires tuning of attraction strength to avoid overly gel-like behavior.

\section{Rubisco scaffold architecture}

\subsection{Rubisco center}

Each Rubisco scaffold contains one center particle of type \(R\). This particle represents the coarse-grained center of the Rubisco scaffold.

\subsection{Patch geometry}

Each Rubisco has eight patch particles of type \(P\). Their initial directions are the normalized cube vertices:
\[
(1,1,1),\; (1,1,-1),\; (1,-1,1),\; (1,-1,-1),
\]
\[
(-1,1,1),\; (-1,1,-1),\; (-1,-1,1),\; (-1,-1,-1).
\]

For each vector \(\mathbf{v}_p\), the unit direction is:
\[
\hat{\mathbf{n}}_p = \frac{\mathbf{v}_p}{\|\mathbf{v}_p\|}.
\]

If the Rubisco center is at position \(\mathbf{R}_r\), the initial position of patch \(p\) is:
\[
\mathbf{P}_{rp} = \mathbf{R}_r + R_{\mathrm{patch}}\hat{\mathbf{n}}_p,
\]
where:
\[
R_{\mathrm{patch}} = 2.5.
\]

\subsection{Patch tether bonds}

Each patch is connected to its Rubisco center by a harmonic tether:
\[
U_{\mathrm{patch}}(r)
=
\frac{1}{2}k_{\mathrm{patch}}(r-r_{\mathrm{patch}})^2.
\]

The default parameters are:
\[
k_{\mathrm{patch}} = 50.0,
\qquad
r_{\mathrm{patch}} = 2.5.
\]

This makes Rubisco a flexible patchy scaffold. The patches are not a mathematically exact rigid body, but they remain localized near the scaffold surface. This implementation is robust and GPU-compatible.

\section{EPYC1 polymer architecture}

Each EPYC1 chain is a 12-bead chain:
\[
[S,T,S,T,S,T,S,T,S,T,S,T].
\]

With zero-based indexing, sticker beads occur at:
\[
\{1,3,5,7,9,11\}.
\]

Adjacent beads are connected by harmonic bonds:
\[
U_{\mathrm{chain}}(r)
=
\frac{1}{2}k_{\mathrm{chain}}(r-r_0)^2.
\]

The default chain parameters are:
\[
k_{\mathrm{chain}} = 10.0,
\qquad
r_0 = 1.2.
\]

The purpose of the chain model is to make EPYC1 a spatially connected multivalent linker rather than six independent binding sites.

\section{Simulation box and initialization}

The simulation uses a cubic periodic box with side length:
\[
L = \texttt{BOX}.
\]

HOOMD boxes are centered at the origin, so valid particle positions lie in:
\[
-\frac{L}{2} \le x,y,z < \frac{L}{2}.
\]

The code uses a centered wrapping function:
\[
\mathrm{wrap}(x) = \left((x+L/2)\bmod L\right)-L/2.
\]

Rubisco centers and EPYC1 chain centers are initialized on grids with small random perturbations. This avoids severe overlaps and numerical blow-ups. Patch particles are not wrapped independently after placement because independent wrapping can create very long patch tether bonds across the periodic boundary.

\section{Forces and potentials}

The model uses:
\begin{enumerate}[label=\arabic*.]
    \item bounded Gaussian soft repulsion,
    \item attractive Gaussian sticker--patch interactions,
    \item harmonic bonds for EPYC1 chains,
    \item harmonic tethers for Rubisco patches.
\end{enumerate}

\subsection{Bounded Gaussian soft repulsion}

The soft repulsion uses a Gaussian potential:
\[
U_{\mathrm{G}}(r)
=
\epsilon \exp\left[-\frac{1}{2}\left(\frac{r}{\sigma}\right)^2\right].
\]

For positive \(\epsilon\), this is a bounded repulsion. It was chosen instead of Lennard-Jones repulsion because LJ repulsion can generate enormous forces if particles overlap. The Gaussian potential remains finite even at \(r=0\).

The default repulsive interactions are:
\begin{longtable}{>{\ttfamily}p{0.22\textwidth} p{0.18\textwidth} p{0.18\textwidth} p{0.24\textwidth}}
\toprule
Pair & epsilon & sigma & cutoff \\
\midrule
R--R & 2.0 & 1.5 & 3.0 \\
S--S & 1.0 & 1.0 & 2.5 \\
T--T & 1.0 & 1.0 & 2.5 \\
S--T & 1.0 & 1.0 & 2.5 \\
R--S & 1.5 & 1.5 & 3.0 \\
R--T & 1.5 & 1.5 & 3.0 \\
P--P & 0.2 & 0.8 & 2.0 \\
P--S & 0.2 & 0.8 & 2.0 \\
P--R & 0.2 & 0.8 & 2.0 \\
\bottomrule
\end{longtable}

All other type pairs are explicitly defined with zero interaction to satisfy HOOMD-blue's type-pair parameter requirements.

\subsection{Sticker--patch attraction}

The EPYC1--Rubisco attractive interaction is between sticker beads \(T\) and patch particles \(P\). It is represented as a Gaussian potential with negative \(\epsilon\):
\[
U_{TP}(r)
=
-\epsilon_{\mathrm{att}}
\exp\left[-\frac{1}{2}\left(\frac{r}{\sigma_{\mathrm{att}}}\right)^2\right].
\]

In the code:
\[
\epsilon_{\mathrm{att}} = |\texttt{KON}|,
\]
so:
\[
\texttt{PATCH\_EPSILON} = -|\texttt{KON}|.
\]

The default attraction parameters are:
\[
\sigma_{\mathrm{att}} = 1.0,
\qquad
r_{\mathrm{cut}} = 3.0.
\]

This interaction is continuous and reversible. Stickers and patches associate when thermal motion and attractive forces bring them together, and they separate when thermal motion overcomes the attraction.

\section{Brownian dynamics}

The model uses HOOMD-blue Brownian dynamics. Conceptually, the update is:
\[
\mathbf{x}(t+\Delta t)
=
\mathbf{x}(t)
+
\mu\mathbf{F}(t)\Delta t
+
\sqrt{2D\Delta t}\bm{\eta}(t),
\]
where:
\begin{itemize}
    \item \(\mathbf{F}\) is the force,
    \item \(D\) is an effective diffusion coefficient,
    \item \(\bm{\eta}\) is Gaussian noise,
    \item \(\mu\) is mobility.
\end{itemize}

The code sets:
\[
k_BT = 1.0,
\]
and uses the following friction coefficients:
\begin{longtable}{>{\ttfamily}p{0.16\textwidth} p{0.20\textwidth} p{0.48\textwidth}}
\toprule
Particle & gamma & Interpretation \\
\midrule
R & 5.0 & Rubisco centers diffuse slowly. \\
P & 5.0 & Patch particles diffuse slowly due to scaffold tethering. \\
S & 1.0 & EPYC1 spacer beads diffuse faster. \\
T & 1.0 & EPYC1 sticker beads diffuse faster. \\
\bottomrule
\end{longtable}

The timestep is:
\[
\Delta t = 0.0005.
\]

\section{Parameter sweep}

The current parameter sets are:
\begin{longtable}{>{\ttfamily}p{0.22\textwidth} p{0.14\textwidth} p{0.14\textwidth} p{0.14\textwidth}}
\toprule
Name & BOX & KON & KOFF \\
\midrule
control & 40 & 1.0 & 0.010 \\
weak & 36 & 2.0 & 0.008 \\
nucleation & 33 & 3.0 & 0.006 \\
early & 31 & 5.0 & 0.004 \\
coarsening & 29 & 8.0 & 0.002 \\
\bottomrule
\end{longtable}

In this HOOMD implementation, \texttt{KON} is used as attraction strength, not as a kinetic on-rate. \texttt{KOFF} is recorded for continuity with earlier workflows, but there is no explicit stochastic off-rate. Unbinding emerges from force-based dynamics.

\section{Density-like quantity}

A useful density-like quantity is:
\[
\eta = \frac{N_R+N_C}{L^3}.
\]

For the default sets:
\begin{align*}
\eta_{\mathrm{control}} &= \frac{72}{40^3} = 0.001125,\\
\eta_{\mathrm{weak}} &= \frac{72}{36^3} \approx 0.001543,\\
\eta_{\mathrm{nucleation}} &= \frac{72}{33^3} \approx 0.002004,\\
\eta_{\mathrm{early}} &= \frac{72}{31^3} \approx 0.002417,\\
\eta_{\mathrm{coarsening}} &= \frac{72}{29^3} \approx 0.002952.
\end{align*}

\section{Contact analysis}

The simulation is force-based, but contacts are detected by distance during analysis. A sticker--patch contact is counted when:
\[
r_{TP} < r_{\mathrm{contact}},
\]
where:
\[
r_{\mathrm{contact}} = 3.0.
\]

The total number of contacts is:
\[
N_{\mathrm{contacts}}(t)
=
\sum_{i\in T}\sum_{j\in P}
\mathbf{1}\left(r_{ij}<r_{\mathrm{contact}}\right).
\]

The CSV writes this quantity as both \texttt{contacts} and \texttt{bonds}. In this model, ``bonds'' means distance-defined sticker--patch contacts, not permanent discrete bonds.

\section{Cluster analysis}

Clusters are computed from a graph whose nodes are:
\begin{itemize}
    \item Rubisco centers,
    \item whole EPYC1 chains.
\end{itemize}

The graph has:
\[
N_{\mathrm{graph}} = N_R + N_C = 72
\]
nodes.

If any sticker bead on EPYC1 chain \(c\) contacts any patch on Rubisco \(r\), the analysis adds an edge between EPYC1 node \(c\) and Rubisco node \(r\).

The outputs are:
\begin{align}
N_{\mathrm{clusters}}(t) &= \text{number of connected components},\\
N_{\max}(t) &= \text{largest connected component size},\\
f_{\max}(t) &= \frac{N_{\max}(t)}{72}.
\end{align}

\section{Outputs}

For each parameter set, the script writes:

\begin{itemize}
    \item \texttt{traj\_patchy\_gpu/<name>.gsd}: native HOOMD trajectory.
    \item \texttt{xyz\_patchy\_gpu/<name>.xyz}: OVITO-readable trajectory.
    \item \texttt{data\_patchy\_gpu/<name>.csv}: time series of contacts, clusters, largest cluster, and largest fraction.
    \item \texttt{plots\_patchy\_gpu/<name>\_bonds.jpeg}: contacts/bonds vs time.
    \item \texttt{plots\_patchy\_gpu/<name>\_clusters.jpeg}: clusters vs time.
    \item \texttt{plots\_patchy\_gpu/<name>\_largest.jpeg}: largest cluster size vs time.
    \item \texttt{plots\_patchy\_gpu/<name>\_fraction.jpeg}: largest cluster fraction vs time.
\end{itemize}

\section{Frame count and meaning}

A frame is one saved snapshot of the system. If:
\[
N_{\mathrm{steps}} = 10000,
\qquad
\texttt{WRITE\_EVERY} = 100,
\]
then one expects approximately:
\[
10000/100 = 100
\]
saved frames, plus any manually written initial frame if included.

Frames determine:
\begin{itemize}
    \item how many timepoints appear in the CSV,
    \item how smooth the plots are,
    \item how many animation frames are available in OVITO.
\end{itemize}

\section{Dual-GPU workflow}

The model uses one HOOMD simulation per GPU process. A dual-GPU sweep runs different parameter sets on different GPUs. A typical launcher is:
\begin{lstlisting}[language=bash]
#!/bin/bash

i=0
for row in $(jq -c '.[]' params.json); do
    if (( i % 2 == 0 )); then
        CUDA_VISIBLE_DEVICES=0 python simulate_hoomd_patchy_full.py <(echo "$row") > gpu0_${i}.log 2>&1 &
    else
        CUDA_VISIBLE_DEVICES=1 python simulate_hoomd_patchy_full.py <(echo "$row") > gpu1_${i}.log 2>&1 &
    fi
    ((i++))
done

wait
\end{lstlisting}

\section{How to reconstruct the model from scratch}

To rebuild the model:
\begin{enumerate}[label=\arabic*.]
    \item Install GPU-enabled HOOMD-blue.
    \item Define four particle types: \(R\), \(P\), \(S\), and \(T\).
    \item Create 24 Rubisco center particles.
    \item Around each Rubisco center, create 8 patch particles in normalized cube-vertex directions at radius 2.5.
    \item Tether each patch to its Rubisco center with a harmonic bond.
    \item Create 48 EPYC1 chains.
    \item Represent each EPYC1 as 12 beads with sequence \([S,T,S,T,S,T,S,T,S,T,S,T]\).
    \item Connect adjacent EPYC1 beads with harmonic bonds.
    \item Add bounded Gaussian soft repulsion between steric particle pairs.
    \item Add attractive Gaussian \(T-P\) interactions.
    \item Integrate using Brownian dynamics.
    \item Save GSD trajectories.
    \item Convert GSD trajectories to XYZ for OVITO.
    \item Detect \(T-P\) contacts by distance cutoff.
    \item Build Rubisco--EPYC1 contact graphs.
    \item Compute clusters, largest cluster, and largest fraction.
    \item Save CSVs and plots.
\end{enumerate}

\section{Current limitations}

The current version is substantially improved, but still has limitations:
\begin{enumerate}[label=\arabic*.]
    \item Rubisco patches are tethered particles, not a mathematically rigid body.
    \item The \(T-P\) attraction is isotropic around the patch particle rather than angularly selective.
    \item \texttt{KOFF} is recorded but not used as a kinetic off-rate.
    \item Contacts are analysis-defined by distance, not explicit chemical bonds.
    \item Parameters are phenomenological and require calibration.
\end{enumerate}

\section{Summary}

The current HOOMD-blue model contains:
\begin{itemize}
    \item 24 Rubisco centers,
    \item 192 Rubisco patch particles,
    \item 48 EPYC1 chains,
    \item 576 EPYC1 beads,
    \item 288 EPYC1 sticker beads,
    \item a 2:1 EPYC1:Rubisco ratio,
    \item a 1.5:1 sticker:patch ratio,
    \item force-based sticker--patch attraction,
    \item explicit EPYC1 polymer connectivity,
    \item explicit Rubisco patch tethering,
    \item Brownian dynamics,
    \item GPU execution,
    \item dual-GPU sweep support,
    \item OVITO, CSV, and plot outputs.
\end{itemize}

The model is designed to explore how multivalent EPYC1 linkers organize patchy Rubisco scaffolds into dynamic clusters or condensate-like assemblies.

\end{document}
